\chapter{Álgebra Conmutativa}

\section{Anillos Conmutativos}

Consideramos $A$ un conjunto, luego un grupo abeliano $(A,+)$ y finalmente añadimos una multiplicación 
\Func{\cdot}{A\times A}{A}{(a, a')}{a \cdot a'}
asociativa, conmutativa y un elemento neutro, $1_A$. Además, se verifica la propiedad distributiva.

\begin{ejemplo}\
    \begin{itemize}
        \item El primer ejemplo que vamos a tener en cuenta es
        \begin{gather*}
            \bb{Z}\subset \bb{Q} \subset \bb{R} = \bar{\bb{Q}} \subset \bb{C} \subset \bb{H} \subset \bb{O}
        \end{gather*}
        Estos son ejemplos clásicos y a partir de ellos podremos construir más.

        \item Consideramos un cuerpo $\bb{K}$ y un conjunto $X\neq \emptyset$. Si tomamos todas las aplicaciones de $X$ en $\bb{K}$ que será un conjunto isomorfo a $\bb{K}^I = \Pi_{i\in I} \bb{K}_i$.\\
        
        Supongamos ahora $X=\bb{N}$ y $Maps(\bb{N}, \bb{K}) \supseteq \bb{K}^{(\bb{N})}$ donde podemos considerar el producto de Hurwitz.
        \begin{gather*}
            (a_0, ..., a_n) \cdot (b_0, ..., b_n) = \sum_{k=0}^{n} \binom{n}{k} a_k \cdot b_{n\cdot k}
        \end{gather*}

        % Si queremos considerar ahora $l\geq 1$ y tomar $X=\bb{N}^l$ 

        Si consideramos $A$ un anillo conmutativo, nos planteamos si podemos hacer esta construcción. Sabemos ya que hay un morfismo inyectivo $\bb{Z} \to A$, por lo que será una $\bb{Z}$-Álgebra.\\

        Sean $A$ y $B$ dos anillos conmutativos. Se dice que $B$ es una $A$-Álgebra si existe un morfismo $A\to B$ de anillos.\\

        % Buscar definición de morfismo de anillos 

        Nos planteamos ahora las siguientes relaciones
        \begin{gather*}
            \xymatrix{
                A[X] & \cong & A \underset{{\bb{K}}}{\otimes}  \bb{Z}[X]\\
                \bb{Z}[X]
            }
        \end{gather*}

        %buscar definición de producto \otimes_{\bb{K}} tensor g

        Si yo cojo $V$ tal que $dim_{\bb{K}}(V)<\infty$ y consideramos su dual
        \begin{gather*}
            V^* = Hom(V, \bb{K})
        \end{gather*}
    \end{itemize}
\end{ejemplo}

% Aqui falta algo (producto tensorial)

% Este producto tensorial desciende a $\bb{K}-$Álgebras (siempre conmutativas)

% ecualizador A\times B \to A \otimes B

Si consideramos $W$ un $\bb{K}-$espacio vectorial, en general
\begin{gather*}
    W \otimes ... \otimes W = W^{\otimes^n}\\
    \prod_{n\in N} W^{\otimes^n}
\end{gather*}
dentro de este último podemos considerar los de soporte finito, que son esas tuplas infinitas, con un número finito de elementos distintos de $0$.

Si $dim(V), dim(W) < \infty$ tendremos que 
\begin{gather*}
    dim(Hom(V,W)) = dim(M_{n\times m}(\bb{K}))
\end{gather*}
lo que nos dice que 
\begin{gather*}
    dim(V^*\otimes W) = n\cdot m
\end{gather*}
queremos introducir ahora las sumas directas, es decir, $\oplus_{n\geq 0} W^{\otimes^n}$ que escribiremos como
\begin{gather*}
    \lambda_0 \overset{\cdot}{+} w_1 + w_2 \otimes w_3 \overset{\cdot}{+} ... \overset{\cdot}{+} w_{k_1} \otimes ... \otimes w_{nt}
\end{gather*}
se usa el producto concatenado para considerar el álgebra.

Si consideramos el ideal $I = \langle \{ v \otimes w + w \otimes v \}_{v\in W,\ u\in W} \rangle$

\begin{gather*}
    \wedge W = T_{\bb{K}} 
\end{gather*}

Antes hemos dicho lo que es un morfismo de anillos $A \overset{\varphi}{\to} B$ que son los que 
\begin{enumerate}
    \item $\ker(\varphi) = \{a \in A : \varphi(a) = 0\}$
    \item $\Im(\varphi) = \{b\in B : \exists a \in A,\ \varphi(a) = b\}$
\end{enumerate}
dado $A$ un anillo, tenemos que $(I,+) \leq (A,+)$. Se dice que $I\leq A$ es un ideal de A por lo que $a\cdot x \in I$ para todo $a\in A$, $x\in I$. De esta forma
\begin{gather*}
    A \cdot I = I \cdot A \subset I\\
    A \overset{\pi}{\to} A/I \text{ en } Ab\\
    a \mapsto a + I = \bar{a}
\end{gather*}
y queremos ver que $\bar{a} \cdot \bar{a'} = \bar{aa'}$ está bien definida, es decir, que no depende de los representantes elegidos. \\

Nos queda ver cuál es la propiedad universal. 

La idea es que si tenemos $X$ un conjunto y $R$ una relación de equivalencia podemos considerar el conjunto cociente
\begin{gather*}
    X/R = \{\bar{x} : x\in X\}
\end{gather*}
donde
\begin{gather*}
    \bar{x} = \{x'\in X : xRx'\}
\end{gather*}
y se dice que $\bar{x} = \bar{y}$ si y solo si $xRy$ donde
\begin{gather*}
    X \to X/R \\
    x \mapsto \bar{x}
\end{gather*}
La propiedad universal nos dice que si tenemos una aplicación $f:X\to Y$ bien definida tal que 
\begin{gather*}
    xRx' \Rightarrow f(x)=f(x')
\end{gather*}
entonces existe una única $\bar{f}$ que hace conmutar el siguiente diagrama
\begin{gather*}
    \xymatrix{
        X \ar[d]_{\pi_R} \ar[r]^f & Y\\
        X/R \ar@{-->}[ur]_{\bar{f}}
    }
\end{gather*}

Con esto podemos hacer m ichas cosas. Por ejemplo si tenemos una torre de ideales $I \leq J \leq A$ podemos considerar los cocientes
\begin{gather*}
    A/I \geq J(I)
\end{gather*}
Si tenemos un cociente, entonces todos los ideales de un cociente. Consideramos $I\leq A$ y el anillo $A/I$, entonces los ideales $k\leq A/I$ considerando la proyección $A \overset{\pi}{\to} A/I$ serán
\begin{gather*}
    J=\pi^{-1}(k) = \{a \in A : a+I \in K\} \leq A
\end{gather*}
ya que para $a\in A, x\in A$
\begin{gather*}
    x \cdot(a+I) = xa + I
\end{gather*}
y lo que queremos ver es que $J/I = K$ y todos los ideales en el cociente se escriben como $J/I$. Además, $J/I$ es un grupo abeliano. Consideramos $x \in J/I$ y $\alpha \in A/I$ y podemos multiplicar
\begin{gather*}
    x\cdot \alpha = \bar{j} \cdot \bar{a} = \bar{ja} = aj + I
\end{gather*}
donde $\alpha = a + I = \bar{a}$. 

Si ahora consideramos 
\begin{gather*}
    \frac{A/I}{J/I} \cong A/J
\end{gather*}
por el teorema de isomorfía.

% No sé lo que son objetos ni flechas

% buscar definición de monomorfismo de a,b,c y morfismo

Sabemos que $\bb{Z}_n$ son todos anillos finitos para $n>1$. Consideramos los retículos de $D(12) = \{1,2,3,4,6,12\}$ y el de los ideales de $\bb{Z}_{12}$ % Hacer los retículos con la notación \langle \bar{1} ¶angle = \bb{Z}_{12}

Vemos así que los ideales de $A$, $Id(A)$ es un retículo, es decir, un conjunto parcialmente ordenado con una operación binaria con ciertas propiedades % buscar propiedades
\begin{gather*}
    I \wedge J = I \cap J\\
    I \vee J = A + J
\end{gather*}
Si consideramos $S \subseteq A$ tenemos 
\begin{gather*}
    \langle S \rangle = \cap_{S \subseteq I \leq A} I = \{\sum_{finita} a_is_i : a_i\in A,\ s_i\in S\}
\end{gather*}
Además,
\begin{gather*}
    \sum_{\lambda \in \Lambda} I_\lambda = \langle \cup_{\lambda \in \Lambda} I_\lambda \rangle
\end{gather*}

\begin{teo}[Teorema de Krull]
    Sea $A$ un anillo conmutativo, $I\leq A$ un ideal propio ($\neq 0, A$), entonces existe $M\in Max(A)$ tal que $I \leq M$
\end{teo}
Siempre se verifica 
\begin{gather*}
    Max(A) \subseteq Spec(A)  \subseteq Irr(A)
\end{gather*}
donde $Spec(A)=\{p\in Id(A) : p \text{ primo }\}$. Diremos que $p\leq A$ es primo si y solo si se verifican 
\begin{enumerate}
    \item $ab\in p \Rightarrow a\in p$ o $b\in p$
    \item $I,J\leq A$ tal que $IJ \subseteq P \Rightarrow I \subseteq I\subseteq P$ o $J\leq P$
    \item $I,J\in A$ tal que $I\cap J \subseteq P \Rightarrow I \subseteq P$ p $J\subseteq P$
\end{enumerate}

La demostración del teorema anterior la haremos con el Lema de Zorn. 
\begin{lema}
    Consideramos $(P, \leq)$ un poset si toda cadena admite un supremo. Entonces $P$ posee un elemento maximal.
\end{lema}

\begin{proof}
    Tenemos $0\lneq I \lneq A$ y consideramos
    \begin{gather*}
        P_I = \{J \leq A : I \leq J\} \neq \emptyset
    \end{gather*}
    y la cadena $\{J_\lambda\}_{\lambda\in \Lambda}$ en $P_I$.
    Considerando $\cup_{\lambda \in I} J_\lambda$ un supremo de la cadena.
\end{proof}

Consideramos $Max(A)\subseteq Spec(A)$, entonces todos los cocientes $A/Q$ donde $Q\in Max(A)$ entonces eso será un cuerpo pero los cocientes $A/P$ donde $P\in Spec(A)$ entonces será un dominio.

\begin{lema}
    Si cogemos $x\in I(A)$, es equivalente a decir que $x\notin M$ para cualquier $M\in Max(A)$.
    \begin{proof}
        \item[$\Rightarrow$)] Trivial
        \item[$\Leftarrow$)] Sabemos que $x\neq 0$ y $x\neq 1$ y suponiendo que $x\notin U(A)$ Krull nos dice que $\langle x \rangle$ es propio y existe un $M\in Max(A)$ tal que $\langle x \rangle \subseteq M$
    \end{proof}
\end{lema}

\section{Anillos locales}

\begin{definicion}
    Si $A$ no nulo se dice que es \textbf{local} si posee un solo ideal maximal, es decir
    \begin{gather*}
        Max(A) = \{x\}
    \end{gather*}
\end{definicion}

\begin{prop}
    La definición anterior es equivalente a decir que $A\setminus U(A)$ es un ideal. En tal caso $M=A\setminus U(A)$ es el único ideal maximal.

    % \begin{proof}
    %     Como $A$ no es nulo Krull nos dice que $Max(A)\neq \emptyset$.
    % \end{proof}
\end{prop}

Consideramos $\bb{K} = \bb{R}$ y $A$ es un $\bb{R}-$Álgebra
% buscar definición de R-Álgebra
y podemos considerar el conjunto
\begin{gather*}
    Alg_{\bb{R}}(A, \bb{R})
\end{gather*}
el conjunto de morfismos de $\bb{R}$-Álgebras. Estas aplicaciones llevan las sumas en sumas, sueltan los escalares y son multiplicativas. Este conjunto es muy importante.
\begin{gather*}
    Alg_{\bb{R}}(A, \bb{R}) \to Max(A)\\
    \varphi \mapsto ker(\varphi)
\end{gather*}
% Teorema de ceros de Hilbert
Si cogemos $A = \frac{\bb{R}[x]}{\langle x^2 + 1 \rangle} \cong \bb{Q}(\sqrt{-1})$, entonces las algebras de $(A, \bb{R})$ es vacío.

Si cojo $X$ un espacio topológico y considero 
\begin{gather*}
    C(x) = \text{funciones continuas de $x$ en $\bb{R}$}
\end{gather*}

Si cogemos un $x\in X$ podemos asociarle un $\mu_x = \{f \in \varphi(x) : f(x)=0\}$ y entonces $X\to Max(\varphi(x))$ es una aplicación bien definida con $x\mapsto \mu_x$.

Que $X$ sea conexo nos dice que $\varphi(x)$ es idempotente.\\

Qué pasa si ahora tomamos un $x_0\in \bb{R}$ y tomamos una función $f:I\to \bb{R}$ tal que $f(x_0)\neq 0$ no podemos tratarlo como un elemento invertible.

La idea es tomar $I$ un abierto de $\bb{R}$ y una función $f:I\to \bb{R}$ y esto nos da un par
\begin{gather*}
    \{(I, f) : I \in V_{x_0}(\bb{R}),\ f\in C(I, \bb{R})\}
\end{gather*}
y vamos a dotarlo de una relación de equivalencia
\begin{gather*}
    (I,f)R(J,g) \sii \exists K\subseteq I\cap J \ \ : \ \ f_{|K} = g_{|K}
\end{gather*}
Llamaremos el germen de $f$ en $x$ a $\bar{(I,f)}$ y tendremos que $C_{x_0}(\bb{R}) = I_{x_0}/R_{x_0}$ es un anillo, y de hecho un $\bb{R}-$Álgebra. Tenemos un anillo y su localización, hemos invertado un formalismo llamado localización que nos permite encontrar un sitio donde es invertible. El germen es invertible, no la función. Stalz de un haz.

Queremos formalizar esto. Para ello tomamos $A$ un anillo y $S \subseteq A$. Decir que $S$ es multiplicativo, es decir
\begin{enumerate}
    \item $0\notin S$
    \item $1\in S$
    \item $s,z\in S \Rightarrow sz\in S$
\end{enumerate}

Por ejemplo, para $A=\bb{Z}$ consideramos $S=\bb{Z}\setminus \{0\}$ y tenemos que $A$ es un dominio. Si cogemos $A\setminus \{0\}$ también es un $S$ multiplicativo.

Si cogemos $P\in Spec(A)$ entonces $S = A\setminus P$.

La idea es coger $S\times A$ (producto cartesiano de toda la vida) y lo dotamos de una relación de equivalencia
\begin{gather*}
    (s,a)R_S(r,b) \overset{def}{\sii} \exists u\in S \ : \ u(ar - sb) = 0
\end{gather*}
% Probar que es una relación de equivalencia
Vamos a denotar la clase de equivalencia de $(s,a)$, es decir $\bar{(S,a)}$ como $\frac{a}{S}$.

Vamos a denotar $\frac{S\times A}{R_S} = AS^{-1}$. 

Nos planteamos cómo definir las operaciones:
\begin{gather*}
    \frac{a}{s} + \frac{b}{r} := \frac{as + sb}{sr}\\\\
    \frac{a}{s} \cdot \frac{b}{r} : = \frac{ab}{sr}
\end{gather*}
y tenemos $0_{AS^{-1}} = \frac{0}{1}$ y $1_{AS^{-1}} = \frac{1}{1}$.
Queremos ver que están bien definidas. 

Hasta ahora hemos hecho un anillo con un conjunto multiplicativo y un morfismo que llamaremos morfismo canónico de la multiplicación
\begin{gather*}
    A \overset{\tau}{AS^{-1}}\\
    a \mapsto \frac{a}{1}
\end{gather*}
Si $A$ es un dominio, entonces $\tau$ es inyectivo.
\begin{gather*}
    \frac{a}{1} = \frac{0}{1} \sii \exists s \in S \ : \ s(a \cdot 1 - 1 \cdot 0) = 0 \sii sa=0 \Rightarrow a=0\ (s\neq 0)
\end{gather*}
Planteamos ahora cuál es la propiedad universal. Si tengo un morfismo de anillos $\varphi: A \to B$ tal que $\varphi(S) \subseteq U(B)$, entonces existe un único $\tilde{\varphi}$ que hace que el siguiente diagrama sea conmutativo
\begin{gather*}
    \xymatrix{
        A \ar[d]_{\tau} \ar[r]^{\varphi} & B\\
        AS^{-1} \ar@{-->}[ru]_{\tilde{\varphi}}
    }
\end{gather*}

Si tenemos el siguiente diagrama
\begin{gather*}
    \xymatrix{
        A \ar[d]_{\tau} \ar[r]^{\varphi} & B\\
        AS^{-1}
    }
\end{gather*}
entonces $\varphi(S)=T$ es un conjunto multiplicativo de B. Entonces podemos encontrar lo que llamaremos \textbf{localizante}
% Diagrama localizante

\begin{gather*}
    \xymatrix{
        A \ar[d]_{\tau} \ar[r]^{\varphi} & B\\
        AS^{-1} & BT^{-1}
    }
\end{gather*}

Si tenemos un morfismo $\varphi: A \to V$ esto induce un morfismo 
\begin{gather*}
    Spec(B) \to Spec(A)\\
    q \mapsto \varphi^{-1}(q)
\end{gather*}

$I \leq A$ esto induce que $I^e \leq AS^{-1}$ , $I^e = \tau(I)AS^{-1} = \langle \tau(I) \rangle$ dentro de $AS^{-1}$.

$J\leq AS^{-1}$ esto induce que $J^{e}=\tau^{-1}(J)$. Tenemos las siguientes propiedades:
\begin{enumerate}
    \item $I^e = \{u\in AS^{-1} \ : \ u = \frac{a}{s} \text{ para algún } a\in I,\ s \in S\}$
    \item Si $I\cap S\neq \emptyset $, entones $I^e = AS^{-1}$
    \item $(J^e)^e = J$ y $I \leq (I^e)^e$
\end{enumerate}

Si $A = \bb{Z}$ localizado en $S$ con $S = \{S^n : n\geq 0\}$.

Si cogemos esto vemos que $\frac{-2}{5} = \frac{-10}{2}$ y sicogemos el ideal $I=10\bb{Z}$ tenemos que $-10\in I$ pero $-2\notin I$.


 



